\section{Smoothness of the Process}
\label{sec:conclusion_smoothness}%
Overall, the installation and validation of the BBP system went smoothly. The READMEs provided clear steps for bringing up the backend and the mobile app; dependency installation completed without errors, and the automated test suites ran successfully after minor environment adjustments (e.g.\ Jest \texttt{testMatch} on Windows, \texttt{NODE\_OPTIONS} in the shell). Manual testing of path creation and trip recording also went well: the app responded as expected in most flows.

\section{Errors and Issues Encountered}
\label{sec:conclusion_errors}%
We did encounter a number of issues, which are detailed in Section~\ref{sec:problems_encountered} . In summary we found:
\begin{itemize}
    \item \textbf{Documentation gaps:} JWT secrets, \texttt{DATABASE\_URL} requirements and details about the need for Redis were not fully spelled out.
    \item \textbf{Platform-specific behaviour:} The backend test script and the app \texttt{testMatch} pattern required workarounds on Windows.
    \item \textbf{Runtime issues:} Keyboard not dismissible in path description; ``Maximum update depth exceeded'' when drawing paths in the emulator; occasional snapping failures; trip recorded on path even when user deviates.
\end{itemize}
None of these blocked acceptance: they either had simple workarounds, were environment-specific, or pointed to refinements rather than fundamental failures.

\section{Assessment of the Product}
\label{sec:conclusion_assessment}%
Despite the issues above, the BBP implementation comes across as a \textbf{professional and functional product}. The backend offers a coherent API, a solid test suite, and a clear separation of concerns. The mobile app provides a modern, responsive interface for searching paths, creating paths, recording trips, and viewing history and statistics. External integrations work in the evaluated scenarios and improve the user experience. The fact that core features work end-to-end and that automated tests pass (with the noted exception) indicates a mature development process and good alignment with the RASD and DD.

\section{Conclusion}
\label{sec:conclusion_final}%
Taking into account the limited time, the evaluated team has delivered a system that is \textbf{fit for acceptance}. Except for the documented errors and gaps, the application behaves correctly, looks professional, and matches the intended functionality described in the reference documents. The issues we found are mostly minor, platform-related, or documentation-related; they do not undermine the conclusion that BBP is a good product that meets the main goals of the project.
